\chapter*{前言}
\addcontentsline{toc}{chapter}{前言}

本书把深度学习写成可运行的对象与映射：用数据构造经验分布，用可微模型给出预测，
用标量目标衡量误差，用优化算法更新参数。笔记只保留定义、成立条件、关键推导和结论。

同一对象用三种等价表示同时出现：概念、公式、可执行计算。
后文以公式为主陈述模型，代码只在需要核验映射时出现。
若把一次观测写成输入 $\bm{x}$ 与目标 $y$，带参数 $\bm{\theta}$ 的模型给出
\[
\hat{y}=f_{\bm{\theta}}(\bm{x}).
\]
$\hat{y}$ 由 $\bm{x}$ 与 $\bm{\theta}$ 算出：回归时它是与 $y$ 同量纲的点预测，不是概率；
分类时它往往先是未归一化分数，再经后续映射才变成 $[0,1]$ 上的类概率。
标量损失
\[
L(\bm{\theta})
=
\frac{1}{n}\sum_{i=1}^{n}
\ell\bigl(f_{\bm{\theta}}(\bm{x}^{(i)}),y^{(i)}\bigr)
\]
把每条样本的预测与真值收成一个非负实数再取平均，表示这批预测整体离真值有多远，
不是某类发生的概率。参数更新
\[
\bm{\theta}
\leftarrow
\bm{\theta}-\eta\,\nabla_{\bm{\theta}}L(\bm{\theta})
\]
沿梯度的反方向走一步：$\nabla_{\bm{\theta}}L$ 的第 $j$ 个分量是 $\theta_j$ 增加 $1$ 个单位时
$L$ 大约增加多少，$\eta>0$ 是步长。
